% ID: SEL-D0082-OB-E04
% ESTADO: cotejada. Segunda pasada visual y borrador compilado revisados.
% Fuente DOC: Selectividad/sel_2012_matematicas_aplicadas/4 MATEMÁT. APLIC. CC. SS/EXAMEN 4-ANDALUCÍA.doc
% SHA-256 original: 8446c5da6b4b7cfc0f48b89d1b2d052ae0c8a1c4b5bc0b86b7e63399ce8656f8
% Página convertida: 2; opción B; ejercicio 4.

Se considera que, a lo sumo, el $5\%$ de los artículos guardados en un almacén son defectuosos. Pasado un tiempo, la persona encargada del mantenimiento del almacén decide investigar si esa estimación es adecuada. Para ello, escoge aleatoriamente 300 artículos de los que 35 están defectuosos.
\begin{enumerate}[label=\alph*)]
\item \textbf{(1.5 puntos)} Plantee un contraste de hipótesis $(H_0:p\leq0.05)$ para determinar si ha aumentado la proporción de artículos defectuosos. Obtenga la región crítica del contraste para un nivel de significación del $5\%$.
\item \textbf{(1 punto)} \textquestiondown{}Qué conclusión se obtiene con los datos muestrales observados?
\end{enumerate}
